% Part: incompleteness
% Chapter: arithmetization-syntax
% Section: substitution

\documentclass[../../../include/open-logic-section]{subfiles}

\begin{document}

\olfileid[fr]{inc}{art}{sub}
\olsection{Substitution}

Rappelons que la substitution consiste à remplacer par un terme~$t$
toutes les occurrences libres d'!!a{variable}~$u$ dans
!!a{formula}~$!A$, à condition que $t$ soit libre pour~$u$ dans~$!A$ ; le résultat s'écrit $\Subst{!A}{t}{u}$.\footnote{Précision éditoriale : cette condition, présente dans la définition antérieure de la substitution, est omise ici dans la source. La récurrence ci-dessous remplace les symboles sans renommer les variables ; hors de ce domaine, elle peut produire une capture.} Cette
opération, effectuée sur les numéros de G\"odel des !!{variable}s,
des !!{formula}s et des termes, est récursive primitive.

\begin{prop}\ollabel{prop:subst-primrec}
Il existe une fonction récursive primitive $\fn{Subst}(x, y, z)$
telle que, lorsque $t$ est libre pour~$u$ dans~$!A$
\[
\fn{Subst}(\Gn{!A}, \Gn{t}, \Gn{u}) = \Gn{\Subst{!A}{t}{u}}.
\]
\end{prop}

\begin{proof}
Nous pouvons définir par récurrence primitive une fonction
$\fn{hSubst}$ comme suit :
\begin{multline*}
\begin{aligned}
\fn{hSubst}(x, y, z, 0) & = \emptyseq \\
\fn{hSubst}(x, y, z, i+1) & =
\end{aligned}\\
\begin{cases}
\fn{hSubst}(x, y, z, i) \concat y & \text{si $\fn{FreeOcc}(x, z, i)$} \\
\fn{append}(\fn{hSubst}(x, y, z, i), (x)_{i}) & \text{sinon.}
\end{cases}
\end{multline*}
Nous définissons alors $\fn{Subst}(x, y, z)$ par
$\fn{hSubst}(x, y, z, \len{x})$.
\end{proof}

\begin{prop}
\ollabel{prop:free-for}
La relation $\fn{FreeFor}(x, y, z)$, vraie si et seulement si le terme
de numéro de G\"odel~$y$ est !!{free for} la variable de numéro de
G\"odel~$z$ dans la formule de numéro de G\"odel~$x$, est récursive
primitive.
\end{prop}

\begin{proof} Exercice. \end{proof}

\begin{prob}
Démontrer la \olref[inc][art][sub]{prop:free-for}.
\end{prob}

\end{document}
